\begingroup
\renewcommand{\arraystretch}{1.45}
\fancyhead[L]{\small ARITHMETISCH-GEOMETRISCHES MITTEL}
\Needspace{43\baselineskip}
\AGMArticle{8}
Wenn wir
\[
1+\tfrac14xx+\tfrac14.\tfrac9{16}x^4
+\tfrac14.\tfrac9{16}.\tfrac{25}{36}x^6+\text{usw.}=y
\]
setzen, so ergibt sich
\[
\begin{aligned}&\tfrac12xx+\tfrac14.\tfrac94x^4\\&\quad+\tfrac14.\tfrac9{16}.\tfrac{25}{6}x^6\\&\quad+\tfrac14.\tfrac9{16}.\tfrac{25}{36}.\tfrac{49}{8}x^8
+\text{usw.}=\frac{x\mathrm{d}y}{\mathrm{d}x}\end{aligned}
\]
und
\[
\begin{aligned}&xx+\tfrac14.9x^4+\tfrac14.\tfrac9{16}.25x^6\\&\quad+\tfrac14.\tfrac9{16}.\tfrac{25}{36}.49x^8+\text{usw.}\\&=\frac{xx\mathrm{dd}y}{\mathrm{d}x^2}+\frac{x\mathrm{d}y}{\mathrm{d}x}\end{aligned}
\]
woraus unmittelbar folgt
\[
\frac{xx\mathrm{dd}y}{\mathrm{d}x^2}+3\frac{x\mathrm{d}y}{\mathrm{d}x}+y
=\frac1{xx}\left(xx\frac{\mathrm{dd}y}{\mathrm{d}x^2}+\frac{x\mathrm{d}y}{\mathrm{d}x}\right)
\]
oder
\[
(x^3-x)\frac{\mathrm{dd}y}{\mathrm{d}x^2}
+(3xx-1)\frac{\mathrm{d}y}{\mathrm{d}x}+xy=0
\]
Auf diese Weise sind also unsere arithmetisch-geometrischen Mittel auf Integralgrößen
zurückgeführt worden und liefern eine partikuläre Lösung dieser Differentialgleichung zweiter Ordnung.

\Needspace{12\baselineskip}
Das vollständige Integral derselben Gleichung ist\hfill\\{}
\(\dfrac{\mathfrak{A}}{\mathrm{M}(1+x,\,1-x)}+\dfrac{\mathfrak{B}}{\mathrm{M}(1,x)}\).

\Needspace{32\baselineskip}
Es sei \(\varphi\) ein beliebiger Winkel. Der Wert des Integrals
\(\int\cos\varphi^2\mathrm{d}\varphi\) von \(\varphi=0\) bis
\(\varphi=\pi\) wird, wie allgemein bekannt ist, \(=\tfrac12\pi\) sein; ebenso ist
der Wert des Integrals \(\int\cos\varphi^4\mathrm{d}\varphi\) zwischen denselben Grenzen
\(=\tfrac12.\tfrac34\pi\); für den Wert des Integrals gilt
\(\int\cos\varphi^6\mathrm{d}\varphi=\tfrac12.\tfrac34.\tfrac56\pi\)
usw. --- schließlich ist, wie von selbst klar ist, \(\int\mathrm{d}\varphi=\pi\).
Hieraus ist ersichtlich, dass der Wert des Integrals
\[
\begin{aligned}&\int\mathrm{d}\varphi\times
\bigl(1+\tfrac12x^2\cos\varphi^2\\&\quad+\tfrac12.\tfrac34x^4\cos\varphi^4\\&\quad+\tfrac12.\tfrac34.\tfrac56x^6\cos\varphi^6+\text{usw.}\bigr)\end{aligned}
\]
oder dieses Integrals \(\int\dfrac{\mathrm{d}\varphi}{\surd(1-xx\cos\varphi^2)}\)
\(=\pi y\) wird, wenn es von \(\varphi=0\) bis \(\varphi=\pi\) genommen
und dabei die Größe \(x\) als konstant betrachtet wird.

\Needspace{25\baselineskip}
Wird nun angenommen, dass die Funktion \(\dfrac1{\surd(1-xx\cos\varphi^2)}\) in eine Reihe der Form
\[
P+2Q\cos2\varphi+2R\cos4\varphi+2S\cos6\varphi+\text{usw.}
\]
entwickelt wird, sodass die Koeffizienten \(P,Q,R,S\) usw. allein von \(x\) abhängen, so wird der vollständige Wert
des oben angegebenen Integrals lauten
\[
P\varphi+Q\sin2\varphi+\tfrac12R\sin4\varphi+\tfrac13S\sin6\varphi
+\text{usw.}+\text{Konst.}
\]
\SourceLeaf{380}
\Needspace{31\baselineskip}
und somit wird der Wert innerhalb der zuvor angegebenen Grenzen $=P\varphi$ sein, woraus $y=P$ folgt. Nun bemerken wir,
dass, weil $\frac{1}{y}=\mathrm{M}(1+x,1-x)$ ist, auch
$\frac{1}{y}=\mathrm{M}(1,\surd(1-xx))$ gilt. Hieraus lässt sich sehr leicht der folgende
allgemeinere Satz herleiten: Wenn ein Ausdruck der Form
$\frac{\alpha}{\surd(\beta-\gamma\cos\varphi^{2})}=W$
in eine nach den Kosinus der Winkel $2\varphi$, $4\varphi$, $6\varphi$ usw.
fortschreitende Reihe entwickelt und deren konstantes Glied mit $P$ bezeichnet wird, die größten
und kleinsten Werte von $W$ aber (nämlich $\frac{\alpha}{\surd(\beta-\gamma)}$ und
$\frac{\alpha}{\surd\beta}$) mit $v,v'$, so wird $\frac{1}{P}$ das arithmetisch-geometrische Mittel zwischen $\frac{1}{v}$ und $\frac{1}{v'}$ sein. Zunächst gelten
die vorangegangenen Überlegungen allerdings nur für den Fall, dass $\gamma$ eine positive
Größe ist; sie lassen sich aber sehr leicht auch auf den Fall ausdehnen, dass $\gamma$ negativ ist.
In diesem Fall ergibt sich nämlich $W=\frac{\alpha}{\surd(\beta-\gamma+\gamma\cos\psi^{2})}$,
wenn $\psi=90-\varphi$ gesetzt wird; somit ist $\frac{1}{P}$ das Mittel zwischen
$\frac{\surd\beta}{\alpha}$ und $\frac{\surd(\beta-\gamma)}{\alpha}$. Wie oben.
\Needspace{28\baselineskip}
Übrigens ist leicht zu erkennen, dass sich derselbe Satz ohne jede Änderung auch auf Ausdrücke der Form $\frac{\alpha}{\surd(\beta+\gamma\cos\varphi)}$ erstreckt, sodass nämlich
$\frac{1}{\text{Konst. Glied}}$ wird
\[
\begin{gathered}\\=\mathrm{M}\left(\frac{1}{\text{Maximalwert}},\frac{1}{\text{Minimalwert}}\right)\\=\mathrm{M}\left(\frac{\surd(\beta-\gamma)}{\alpha},
                 \frac{\surd(\beta+\gamma)}{\alpha}\right)\end{gathered}
\]
denn solche Funktionen werden auf die Form
$\frac{\alpha}{\surd(\beta-\gamma+2\gamma\cos\psi^{2})}$ zurückgeführt, die der eben besprochenen
völlig ähnlich ist, wenn $\varphi=2\psi$ geschrieben wird\HistoricalNote{29}.

\Needspace{24\baselineskip}
Schließlich bemerken wir, dass im Folgenden ein viel allgemeinerer Beweis derselben
Sätze aus ursprünglicheren Prinzipien gegeben werden wird; außerdem werden wir bald auch lehren, alle
übrigen Koeffizienten $Q,R,S$ usw. durch ebenso bequeme Verfahren zu ermitteln.
An dieser Stelle wird diese Untersuchung auch den Nutzen haben, diese Forschungen denjenigen
lieber zu machen, die für die Erhabenheit und göttliche Schönheit ewiger Wahrheiten keinen Sinn haben
und ihren Wert allein nach dem Nutzen zu beurteilen wissen, den sie den Zweigen der
angewandten Mathematik bringen können. Niemand verkennt, von welch großem Nutzen
eine Entwicklung der Koeffizienten $P,Q,R,S$ usw., die so rasch ist wie die aus diesen Prinzipien
hervorgehende, für die physische Astronomie oder die Theorie der Planetenstörungen ist.


\par\begin{center}\rule{.15\linewidth}{.4pt}\end{center}
\clearpage
\endgroup
